% Section 1 of "Communication, Coherence and Collective Agents".
% src: F_brief (purpose, the author's four ideas, H1 to H4, A5 to A11, findings of 2026-11-07); M_spec Sections 0, 2 (decisions 1 to 3, 6, 11, 12, 16, 17), 6; E_edits_digest (main results, statuses and open problems of every section); H_brief Section 0 and M_spec_H5 Sections 0, 2, 4 (the fifth idea, H5, decisions 19 to 22); sections sec:met, sec:hie, sec:ren, sec:env and revised sec:col in final form; E3 proof checks (statuses of the H5 results)
\section{Introduction}\label{sec:int}

Many systems that act in the world are made of smaller systems that act. Cells coordinate to build tissues, animals move in flocks and schools, and robot swarms complete tasks that no single robot can. Other collectives have no single body: cultures, economies, nations and cities are held together by conventions, exchange, law and co-presence, and their members communicate through languages, prices, statutes, archives and meetings. In both kinds of case it is natural to speak of the collective as an agent that senses and acts as a unit. This manuscript asks how communicating agents that perform active inference \citep{friston2017active,parr2022active} form such collective agents, which properties of their communication decide where collectives form and how large they can become, and when a collective is itself an agent with an environment of its own.

Active inference describes a collective as a blanket of blankets, a set of agents whose joint states are screened from the rest of the world by a Markov blanket assembled from the members' own blankets \citep{kirchhoff2018markov,ramstead2018answering,palacios2020markov}, and active-inference models of collectives range from morphogenesis \citep{kuchling2020morphogenesis} to collective behavior derived from surprise minimization \citep{heins2024collective}. In the social sciences, \citet{deutsch1953nationalism} defined a people by the complementarity of its members' communication, and \citet{hayek1945use} held that prices communicate dispersed knowledge. As far as we know, neither line derives a blanket from communication between inferring agents, and neither says what selects a collective among the many sets of agents that have blankets or what limits its size.

We stay within traditional active inference. Every generative model is a model of external physical states, the same world for every agent, and agents interact only through their blankets: a direct message, a persistent record and a performative act are all paths from one agent's active states through external states to another agent's sensory states (\cref{fig:fw-channels,prop:fw-transfer-kernel}). Frames, the private coordinates in which an agent may store its beliefs, are optional encodings whose effects appear only through communication. Most results are proved in the linear-Gaussian reference case, with world $\eta_{t+1}=F\eta_t+w_t$, $w_t\sim\Ncal(0,Q)$, and Gaussian beliefs $q_i=\Ncal(m_i,P_i)$ about $\eta_t$.

\subsection{Five ideas}\label{sec:int-ideas}

The framework develops five ideas of the author, which we state first in his terms. First, meta-agents form locally, with Markov blankets of active, sensory and internal states of their own. Second, communication is an action and receiving a message is a sensory process: agent $j$ projects its internal belief to a subspace through a map $T$, and agent $i$ lifts the message into its full internal belief space through a map $O$. Third, the coherence of message passing is part of the mechanism by which meta-agents form. Fourth, the informational correlation length and correlation time grow as the scale of meta-agents grows: media that carry messages farther, such as the telephone, or keep them longer, such as archives and international organizations, support larger and longer-lived collectives. Fifth, a meta-agent should, or might, be a bona fide agent with internal and external states of its own, and its external states would be a sort of abstract physics for it: a corporation need not sense a physical environment but a more abstract informational one, people sense non-physical things such as the economy, and nerve cells sense chemical potentials but not gravity. Levels are linked as hierarchical active inference links them in organisms, by top-down priors and bottom-up evidence, so that a person receives information from the meta-agents a level above, a culture, a job or a community, and the passage between levels is related to the work of Friston and colleagues on renormalization.

The second idea has an exact form (\cref{sec:fw-codes,sec:fw-reception}). A modality $c$ has a public code $T_c$ of rank $r_c$, and agent $j$ may store its belief in private coordinates $\mu_j=U_jm_j$, so the message that $i$ receives from $j$ is $y_{ij}=T_cU_jm_j+\nu_{ij}=H_jm_j+\nu_{ij}$ with $\nu_{ij}\sim\Ncal(0,R_c)$; double indices are receiver first. The map $T$ is the effective code $H_j=T_cU_j$, which keeps the sender's mean only modulo $\ker H_j$. The map $O$ is the right inverse $O_{ij}$ of the receiver's code model that places a coded value at the belief nearest to the receiver's prior mean in the receiver's own metric; on a noisy channel the receiver first decodes, and the composite is the Kalman gain. With a correctly read noiseless code, $O_{ij}H_j$ is a projector, orthogonal in the receiver's metric, onto a subspace that the receiver's prior selects, and the divergence that the code fails to transmit has an exact formula (\cref{prop:fw-prior-completion}). Understanding is prior-dependent completion. The first, third and fourth ideas need a definition of a collective that a set of agents can fail, measures of coherence that agreement alone cannot satisfy, and lengths defined on the law that communication produces, not the shared world; the fifth needs external states of a meta-agent that are declared before the analysis, and a criterion for when they have a law of their own.

\subsection{Hypotheses and standing assumptions}\label{sec:int-hypotheses}

We state the hypotheses in the final forms that the following sections establish. Secondary statements, eight extensions that address open problems posed in the literature, the proofs and the numerical checks are in the supplement, which \cref{sec:int-organization} describes and whose numbers carry the prefix S. Four objects recur in the hypotheses. The world-conditional law of the agents' errors $e_i=m_i-\eta$ is their law given the whole world path and the channel noise shared by several receivers (\cref{def:col-world-conditional}); conditioning removes the common mode that a shared world induces whether or not agents communicate (\cref{lem:fw-common-mode}). The moment bound (M) asks that this law be stationary with finite second moments (\cref{asm:sca-standing}). A precision is diffusive when it is a sum of nonnegative leaks and a graph Laplacian with nonnegative couplings: each agent's error is pulled toward those of the agents it is coupled to, and toward zero by its leak, which collects grounding and forgetting (\cref{sec:col-world}). The contrast $r(C)=\xin(C)/\xout(C)$ of a region $C$ is the slowest relaxation rate of disagreement inside $C$ divided by the relaxation rate of the members' average, both computed from the diffusive world-conditional precision (\cref{def:col-contrast}).

\begin{hypothesis}[Collectives as graded meta-agents]\label{hyp:int-h1}
Let roles be assigned to channels: a member's channel that reads a non-member or the world is sensory, a member's channel that a non-member reads is active, and every other variable of the members is internal (\cref{def:col-partition}).
\begin{enumerate}
\item[(a)] (Boundary.) With per-channel noise, the canonical partition of a set of agents is an exact path blanket from every conditionally independent start, interior and exterior trajectories being conditionally independent given the boundary trajectories, if, and generically only if, no sensory channel superposes an interior and an exterior source; a region with finite interior whose shell is thick on the scale of an explicit bound on the Markov length $\xi_M$ is an approximate stationary blanket of the world-conditional law when that law's covariance decays exponentially and is bounded with bounded inverse on a graph of polynomial growth, and the correlation length does not control $\xi_M$. The boundary property holds for every candidate of these classes and selects none.
\item[(b)] (Contrast.) Under the moment bound and with a diffusive world-conditional precision, a collective at level $r_0$ is a region $C$ with $r(C)\ge r_0>\max\{\varepsilon^{\max}_C/\bar\varepsilon_C,\bar r_{\mathrm{hom}}\}$, where $r_0$ is declared in advance, $\varepsilon^{\max}_C/\bar\varepsilon_C$ is the ratio of the largest to the mean precision leak on $C$, and the homogeneous baseline $\bar r_{\mathrm{hom}}$ is the supremum of contrast over regions of every size in a comparison medium declared with the model, which keeps the population's declared geometry and replaces its couplings and leaks by their means, or, for a population without declared geometry whose couplings a declared mechanism generates, is the exchangeable state of that mechanism (\cref{def:col-baseline}), the whole population included when it is finite and closed. The ratio $r(C)/\bar r_{\mathrm{hom}}$ is the significance of $C$.
\item[(c)] (Closure, interaction necessity and agency.) A collective is a meta-agent to the degree that a declared coarse variable of its internal states closes with small memory; its structure differs from the applicable declared nulls that delete, freeze or randomize communication or remove selection by expected free energy, and survives the replacement of literal reading by exact fusion; and its members' joint policy matches a group posterior, as measured by an agency index.
\end{enumerate}
\end{hypothesis}
% src: F_brief H1; M_spec decisions 1 to 4, 11, 16; M_spec_H5 decision 22 (part (b), comparison medium with declared geometry); thm:col-boundary, def:col-collective, def:col-baseline, def:agy-meta-agent, thm:agy-h1

Hypothesis H1 in this form is assembled in \cref{thm:col-boundary,thm:agy-h1} for linear-Gaussian agents with fixed or exogenously scheduled attention: part (a) is a theorem, and blanket existence selects nothing; part (b) is a definition, whose baseline is finite on homogeneous lattices (\cref{thm:col-lattice-contrast}) and whose diffusive precision is proved for two classes of dynamics but not in general (\cref{thm:col-diffusive-precision}); and in part (c) the bound on the memory of the collective average, the vanishing of the interaction order parameter under the deletion nulls and the group posterior of a stated deliberation protocol are proved under the conditions of \cref{thm:agy-closure-average,prop:agy-nulls-vanish,thm:agy-agency}, while agency does not follow from contrast, and differing from partner randomization and surviving exact fusion are requirements of \cref{def:agy-meta-agent} (Supplement, \cref{sup:int-status-of-hypothesis-h1}).

Receivers read messages at one of three interpretation levels (\cref{def:fw-levels}). At level 0 a message is an observation of $\eta$ with an assigned noise; at level 1 it is an estimate of $\eta$ whose error, with a covariance learned or read from a law header posted by the sender, is taken to be independent of everything the receiver knows; exact fusion uses the law of the message given the world and the receiver's retained information under the true joint process. Levels 0 and 1 are literal.

\begin{hypothesis}[Communication as action, reception as sensing]\label{hyp:int-h2}
Speaking is an active state that writes a coded function of the sender's belief into an external channel, and receiving is a Bayesian update with a likelihood for the message, which lifts the message into the receiver's full belief space by prior completion. The sender's belief shares information with the receiver's, through a common prior, shared sources and the world's own innovations, and it carries the sender's prior in directions the sender never observed. A literal receiver, which reads a message as a measurement of the coded coordinates of $\eta$, therefore misreads its joint structure: it overvalues messages whose errors are redundant with its own, which is double counting, and undervalues messages whose errors complement its own or whose content completes what it already knows. Exact fusion, which reads messages under the true joint law, is the calibrated reference, and it needs provenance.
\end{hypothesis}
% src: F_brief H2; M_spec decision 6; sec:mis summary

Hypothesis H2 holds in stated scopes (\cref{sec:mis}): overvaluation is proved in the scalar track model, for a band of positive error covariances, and in static common-source worlds that are scalar, aligned or read through invertible noiseless codes; undervaluation is proved for complementary errors in the track model and by closed-form examples, and is elsewhere found by numerical search in two random ensembles with one seed, with exact conditions open (Supplement, \cref{sup:int-status-of-hypothesis-h2}).

\begin{hypothesis}[Coherence and formation]\label{hyp:int-h3}
The coherence of message passing has three components, each scored jointly with accuracy, since coherent but inaccurate message passing is an echo chamber: fidelity, the information about external states that a message carries; calibration, the agreement between the precision a receiver assigns to a message and the message's actual reliability; and consistency, the composition of the translations that receivers apply around cycles. Collective interiors are regions of high fidelity, calibrated precision and consistent translation, and boundaries lie where calibrated precision is low and translations are inconsistent. As far as it is established, this formation clause holds in the following form. In a two-group model at infinite policy precision with exact in-group translations, where the declared start is itself the split, the chain from calibrated precision through attention and couplings to collectives, and convergence to them, are theorems, the collectives at the levels of the window of \cref{thm:frm-split-contrast}. At finite policy precision the theorems are reachability of the split branch above a threshold $\gamma_S$ and, above a threshold $\gamma_T$, a single fixed point on that branch that attracts every state; both are about the frozen-calibration mean-field map, with the stochastic-approximation step from the full model to the map open (\cref{open:frm-closure}) and the contrast of the reached state computed at one parameter point. Formation there is a calibration-level effect of misreading joint structure: receivers that know their translation errors pay no mistranslation cost, so from a start law exchangeable given the world path the model's collectives fail the translation-aware form of the exact-fusion persistence clause, which is proved (\cref{prop:frm-exact-fusion}) and leaves the model without a formation result in the sense of \cref{def:frm-target}, while for full cross-covariance-aware exact fusion this conclusion is heuristic; above the threshold $\gamma_T$ they coincide with the code partition. Persistence under exact fusion with drifting frames is a conjecture (Supplement, \cref{conj:frm-drifting-frames}). Online calibration from a symmetric start reaches the attention pattern of the split into the two groups (not necessarily a split selected by contrast) in 0 to 7 of 8 replicates per cell, with reruns of a cell differing by up to three replicates, and is open (\cref{obs:frm-online}; Supplement, \cref{open:frm-online}). Formation beyond this model is a conjecture (\cref{conj:frm-formation}).
\end{hypothesis}
% src: F_brief H3; M_spec decisions 7, 9, 12; def:coh-components, sec:coh-formation, sec:frm-status; formation clause in the canonical wording of sec:frm-status

The components are developed in \cref{sec:coh} and the two-group model in \cref{sec:frm}; the status of the formation clause is the one stated in the hypothesis, and the theorems behind the components, on withdrawal, holonomy, amplification around cycles and the grounding tension (\cref{thm:coh-withdrawal,thm:coh-holonomy,thm:coh-amplification,thm:coh-grounding-tension}), and behind the formation model (\cref{thm:frm-formation,thm:frm-split-contrast,prop:frm-nulls,prop:frm-exact-fusion}), are summarized with their scopes in the supplement (\cref{sup:int-status-of-hypothesis-h3}).

\begin{hypothesis}[Scale]\label{hyp:int-h4}
Let the world-conditional errors of a network with a schedule independent of the noises, local passing and local bounded noise, (Lin), (Loc) and (B) of \cref{asm:sca-standing}, satisfy the moment bound (M) and a named contraction of the error dynamics: row sums, row and column sums, a spectral bound, domination by a reversible chain, or symmetric translation-invariant passing. Then the informational correlation length grows like $\ell\sqrt{a/\varepsilon}$ under nearest-neighbor passing on lattices, with hop length $\ell$, message weight $a$ and leak $\varepsilon$, and like the square root of the second moment of the couplings when links are longer; the correlation time of agents and media together is $\tau=-1/\ln\rho$, with $\rho$ the spectral radius of their joint transition matrix, which is the inverse spectral gap to leading order as $\rho\to1$ and which records lengthen exactly under an explicit criterion; and, for a diffusive world-conditional precision with uniform leak, the largest extent of a collective at level $r_0$ obeys $\Lstar\le\xi_2/\sqrt{r_0-1}$, where $\xi_2$ is the second-moment length of that precision. Apart from the exact Ornstein--Zernike and second-moment lengths and the exact secular root and memory criterion of the record model, these are upper bounds. Without (M) only finite-time light-cone statements hold, and an amplifying profile is by definition a rumor, not a collective.
\end{hypothesis}
% src: F_brief H4; M_spec decisions 3, 13, 16; thm:sca-h4, cor:col-lstar, rem:sca-rumors

Hypothesis H4 in this form is \cref{thm:sca-h4}: light cones need only locality (\cref{thm:sca-lightcone}), every length, time and volume needs (M), which the agents' own attention can violate, and none of the bounds says that collectives of the bounded size form (Supplement, \cref{sup:int-status-of-hypothesis-h4}).

The framework also makes seven standing assumptions explicit, numbered after the first four hypotheses.

\setcounter{premise}{4}
\begin{premise}[Individualism]\label{asm:int-a5}
Each agent selects its policies by its own expected free energy, and there is no team objective: collectives are attractors of individual expected-free-energy minimization.
\end{premise}
\begin{premise}[Timescales]\label{asm:int-a6}
Inference over rounds is fastest, learning of codes and partner models is slower, attention habits are slower still, and model structure and frames are slowest or fixed.
\end{premise}
\begin{premise}[Recursion with memory]\label{asm:int-a7}
Coarse-graining generically produces memory, so the blanket and the model of a collective may need memory variables.
\end{premise}
\begin{premise}[Abstract locality]\label{asm:int-a8}
Distance lives on the communication graph or on a space of roles, not necessarily in physical space.
\end{premise}
\begin{premise}[Multiplexity]\label{asm:int-a9}
Language, exchange and law are separate modalities, each with its own correlation length; collectives nest within a modality and can overlap across modalities.
\end{premise}
\begin{premise}[Identity under turnover]\label{asm:int-a10}
A collective is a persistent pattern in coarse-grained states, not a membership list.
\end{premise}
\begin{premise}[The grounding tension]\label{asm:int-a11}
The leak $\varepsilon$ caps the correlation length but makes the collective informative about the world; pure listening maximizes reach and produces echo chambers.
\end{premise}
% src: F_brief A5 to A11

Several assumptions become results: \cref{asm:int-a7} becomes the closure identity (\cref{thm:agy-closure-identity}), \cref{asm:int-a10} the identity time under turnover (\cref{prop:agy-identity}), \cref{asm:int-a11} the grounding tension (\cref{thm:coh-grounding-tension}) and \cref{asm:int-a9} the multiplex populations of \cref{sec:abs}; under \cref{asm:int-a8} the size bound holds for every metric on a space of roles (\cref{thm:col-contrast-size}), no theorem takes \cref{asm:int-a6} as a hypothesis, and the theorems rest on further hypotheses named in their statements (Supplement, \cref{sup:int-assumptions-that-become-results}).

The fifth idea is developed in \cref{sec:met,sec:hie,sec:ren,sec:env}, which leave the selection of collectives to H1 and ask whether sets of agents are agents in the sense of \cref{sec:fw}, how levels pass priors and evidence, what coarse-graining preserves and what a sensor optimal for a declared predictive objective reads at each level. Three further objects recur there. The coarse internal, sensory, active and external states $y_C$, $s_C$, $a_C$ and $\eta_C$ of a set $C$ of agents are linear functions of the joint process of world, records, channels and internal states, declared before the analysis, and $b_C=(s_C,a_C)$ is the coarse blanket (\cref{def:met-coarse}). A coarse variable is exactly closed when the kernel of its declared map is invariant under the dynamics, so that it obeys a law of its own given its input, closed of order $d$ when $d$ adjoined linear coordinates make it exactly closed, and approximately closed when a stated memory measure lies below a stated bound; the abstract physics of a meta-agent is the exactly or approximately closed law of its coarse exterior given its coarse act (\cref{def:met-abstract-physics}). The coarse medium of a partition is the Galerkin precision of its block indicators, the precision of the block values when every block is held in internal agreement, and the collectives of level $l$ of a hierarchy are judged by their coarse contrast in the coarse medium of level $l-1$ against that medium's own comparison medium (\cref{def:col-coarse-medium,lem:col-coarse-medium,def:col-hierarchy}). These sections rest on one standing assumption that A5 to A11 do not contain, numbered after them.

\begin{premise}[Declared coarse-graining]\label{asm:int-a12}
The coarse states of a meta-agent are linear functions of the joint process, declared before the analysis together with the partitions of a hierarchy and the geometry of every comparison medium. Whether a declared coarse variable has a law of its own is a question about closure, answered by theorems and never assumed.
\end{premise}
% src: M_spec_H5 decisions 19 (coarse external states declared as linear maps; abstract physics a theorem-level question about closure) and 22 (comparison medium declared with its geometry, hierarchies judged level by level); def:met-coarse, def:col-baseline, def:col-hierarchy

\begin{hypothesis}[Meta-agents as agents, across levels]\label{hyp:int-h5}
Let the coarse states of a set $C$ of agents be declared as in \cref{asm:int-a12}.
\begin{enumerate}
\item[(a)] (Meta-agents are agents.) The coarse exterior has an exact abstract physics if and only if the kernel of its map is invariant under the exterior dynamics and the coarse exterior reads the members' channels only through the coarse act; for block sums the kernel condition is lumpability, and every declared coarse exterior has a unique smallest closed completion. If, in a stable joint process with per-channel noise, the coarse exterior has an exact abstract physics with a detectable coarse filter, the coarse sensory state is the information-form statistic of memoryless readers of the coarse exterior, the exterior reads the members only through the coarse act, the members pool information contributions by sum-preserving consensus, and the remaining hypotheses of \cref{thm:met-h5a} hold, then $b_C$ is an exact path blanket from every conditionally independent start, the coarse internal state is exactly closed given $b_C$, the members' average estimate is, up to the response to their averaged computation noise, the predictive mean of the steady-state coarse filter at every number of consensus rounds, which with the steady-state covariance parametrizes the meta-agent's predictive law of its current exterior from the centralized start with stationary initial covariance, and, under the deliberation protocol of the group model with a shared preference precision, the joint act follows, in the stationary law of deliberation at a decision state, a group posterior with habit tilts bounded by the members' disagreement. These properties describe a meta-agent and select none.
\item[(b)] (Hierarchy.) A member's two-level model of its collective as a latent cause is exact if and only if the collective coordinate is closed. Exact inference on trees sends likelihoods, equivalently precision-weighted prediction errors, up and leave-one-out priors down. A receiver that reads a message containing its own evidence as an independent prior double counts: in a scalar tree one reading leaves it overconfident by a factor below two, and a record rebuilt every round from its readers' posteriors becomes overconfident without bound unless it keeps each contribution once. A record that members write and read is a top-down channel; without decay it lengthens the memory of the collective mode by a factor of about $1+\beta/\omega$, and for grounded coordinates it raises a shared-error floor that the members cannot see in their social and top-down innovations. Linear conformity to a self-constituted prior selects no partition that the declared couplings and readership do not contain.
\item[(c)] (Abstract environments and selective sensing.) Linear sensors with a fixed number of channels or a Shannon budget that carry the most information about a linear-Gaussian exterior at a horizon read the leading canonical variates between present and future: as the horizon grows, the same variates in reversible exteriors, fewer of them under a budget, asymptotically the slowest mode amplitudes in diagonalizable non-reversible ones, and coarse spatial averages in diffusive ones. An added channel has value only if it is correlated with the exterior at the agent's horizon given what the agent already reads, and pragmatic value only if it bears on outcomes that the agent's acts can compensate. Type averages of many agents are, together with the world, exactly closed external states when the agents form interchangeable types, and the average of uncoupled first-order autoregressive agents carries memory when their persistences differ, with long memory in the infinite-population limit under a common shock alone when the density of the persistences $\alpha$ diverges at the unit root like $(1-\alpha)^{\mathfrak p-1}$ with $\frac12<\mathfrak p<1$.
\item[(d)] (Renormalization.) Decimation and block averaging map the world-conditional Gaussian law exactly; the axis rate and the second-moment length flow exactly, nearest-neighbor precisions on $\Z$ and the nested medium renormalize in closed form, and block averages in general leave the diffusive class, so collectives of collectives are judged in coarse media. The static error and the memory of a reduction to block coordinates are one matrix, small when the smallest rate of the eliminated internal modes is large against the equitability defects of the blocks; the elimination of internal states in an adiabatic reduction drops an exact memory kernel; the size bound caps the depth of a hierarchy on a lattice; and in the nested medium without declared geometry, collectives exist at every level only if the ratio $\mathfrak q$ of an agent's total couplings at successive levels is below $b/(2b^2-1)$, with $b$ the branching of the hierarchy, and at no level if it is at least $1/(2b-1)$.
\end{enumerate}
\end{hypothesis}
% src: H_brief Section 3 (H5a to H5d); M_spec_H5 Section 0 and decisions 19 to 22; thm:met-abstract-physics, thm:met-h5a, rem:met-no-selection; prop:hie-closure-exact, thm:hie-messages, prop:hie-literal-once, prop:hie-literal-records, prop:hie-record-gain, thm:hie-floor, lem:hie-record-blind, prop:hie-conformity; thm:env-canonical, prop:env-horizon, prop:env-diffusive, prop:env-selection, prop:env-aggregate-closure, prop:env-aggregate-memory; def:ren-maps, thm:ren-flows, thm:ren-decimation, thm:ren-nested, ex:ren-block-average, lem:ren-galerkin, thm:ren-memory, prop:ren-adiabatic, prop:ren-depth, prop:ren-per-level

Hypothesis H5 in this form is proved in linear-Gaussian models, with exceptions named in the supplement (\cref{sup:int-status-of-hypothesis-h5}): part (a) is \cref{thm:met-h5a}, an exact conditional theorem whose approximate form is open beyond a conditional special case (Supplement, \cref{open:met-approximate}) and which is not a criterion for which sets of agents are meta-agents (\cref{rem:met-no-selection}); parts (b), (c) and (d) collect statements of \cref{sec:hie,sec:env,sec:ren} proved in the models named there, formation by conformity being a conjecture, the depth bound for decimated media in two and more dimensions resting on an inequality proved in one dimension and checked for a single decimation step in two, and the first-order rate of the error of the memoryless reduction a numerical observation.

\subsection{Main results}\label{sec:int-results}

\Cref{tab:int-results-a} lists the main results of \cref{sec:fw,sec:mis,sec:col,sec:agy} with their status. A theorem, proposition, lemma or corollary has a written proof in the text or in the proofs part of the supplement. An example is a closed-form instance or a counterexample. A numerical observation is established only by the computation it reports, and conjectures and open problems are labeled as such. The last column gives the model or the hypotheses in which a result holds, and no result is claimed outside them. Many of the proved results were also checked numerically on small instances (Supplement, \cref{app:checks}).

\begingroup
\small
\setlength{\tabcolsep}{3pt}
\begin{longtable}{@{}>{\raggedright\arraybackslash}p{0.30\textwidth}>{\raggedright\arraybackslash}p{0.17\textwidth}>{\raggedright\arraybackslash}p{0.14\textwidth}>{\raggedright\arraybackslash}p{0.33\textwidth}@{}}
\caption{Main results of \cref{sec:fw,sec:mis,sec:col,sec:agy}.}\label{tab:int-results-a}\\
\toprule
Result & Where & Status & Model or hypotheses\\
\midrule
\endfirsthead
\toprule
Result & Where & Status & Model or hypotheses\\
\midrule
\endhead
\bottomrule
\endfoot
Prior completion: $O_{ij}H_j$ is a projector in the receiver's metric & \cref{prop:fw-prior-completion} & proposition & Gaussian receivers, correct noiseless code; meaning lost for any prior\\
When each interpretation level is exact & \cref{prop:fw-level-optimality} & proposition & belief messages; level 1 fails under a proper common prior\\
Common mode and connected part of the error & \cref{lem:fw-common-mode} & lemma & level-0 network, $\rho(W)<1$; mistranslation in a stationary world\\
\addlinespace[2pt]
Overvaluation of redundant messages (double counting) & \cref{prop:mis-overvaluation-scope}; Supplement, \cref{prop:mis-track-band} & proposition & scalar, aligned or noiseless invertible-code static worlds; track model, positive band\\
Believed-to-true ratios diverge like $(1-\phi)^{-2}$ & \cref{prop:mis-closed-ratios} & proposition & symmetric scalar common-source world with channel noise\\
Undervaluation beyond the proved examples & Supplement, \cref{obs:mis-random-codes} & numerical observation & two ensembles of random static worlds, one seed; exact conditions open\\
Echo chambers exist; exact inference forms no weak or self-sealed one & \cref{thm:mis-echo-chambers,prop:mis-exact-no-echo} & theorem; proposition & closed groups, homogeneous grounding, reversible weights; exact inference with fixed schedules\\
\addlinespace[2pt]
Path blankets characterized; blanket existence selects nothing & \cref{thm:col-path-blanket}; Supplement, \cref{prop:col-selects-nothing} & theorem; proposition & per-channel noise, conditionally independent start\\
$\xi_M$ and $\xi$ are logically independent; explicit bound on $\xi_M$ & \cref{prop:col-markov-vs-correlation}; Supplement, \cref{thm:col-markov-bound} & proposition; theorem & explicit families; well-conditioned decaying covariance\\
Diffusive world-conditional precisions, and counterexamples & \cref{thm:col-diffusive-precision}; Supplement, \cref{ex:col-non-diffusive} & theorem; example & nonnegative weights, reversible or normal with isotropic noise; can fail otherwise\\
Homogeneous lattices have bounded contrast & \cref{thm:col-lattice-contrast} & theorem & $\Z^D$, translation-invariant symmetric kernels positive on axis steps, finite second moments, uniform leak\\
Size bound $\Lstar\le\xi_2/\sqrt{r_0-1}$, sharp up to $\pi/\sqrt{12}$ & \cref{thm:col-contrast-size,cor:col-lstar} & theorem; corollary & diffusive precision, any coupling range, uniform leak\\
Bridge to correlation criteria & \cref{thm:col-bridge} & theorem & uniform leak; brackets for two equal blocks forming the population\\
A finite closed population is a collective of its own comparison medium only if more connected than that medium & Supplement, \cref{prop:col-whole-population} & proposition & uniform leak; never under the mean-coupling exchangeable clause, on tori or on vertex- and edge-transitive declared graphs\\
$\xi_M\le B(1+\xi)$ without conditioning & Supplement, \cref{conj:col-markov-length} & conjecture & networks on $\Z^D$, local noise, row-sum contraction\\
\addlinespace[2pt]
Exact closure if and only if retained agents hear no eliminated one & \cref{thm:agy-closure-identity} & theorem & linear recursions, arbitrary real matrices; coordinate selections\\
Contrast bounds the memory of the collective average & \cref{thm:agy-closure-average} & theorem & reversible diffusion with drift the precision block of $C$, exterior fixed; for a general drift a drift-based contrast, whose relation to $r(C)$ is open\\
Consensus on contributions approaches the union posterior geometrically & \cref{prop:agy-consensus} & proposition & symmetric doubly stochastic consensus; static and dynamic worlds\\
Agency: a group posterior when members agree on redundancy and a belief carries it & \cref{thm:agy-agency,thm:agy-heterogeneous} & theorem & deliberation protocol; explicit distances otherwise\\
Hypothesis H1 in final form & \cref{thm:agy-h1} & theorem & linear-Gaussian, fixed or exogenous attention, (M)\\
\end{longtable}
\endgroup
% src: E_edits_digest main results of fw, mis, col, agy; statuses confirmed in 02_framework.tex, 03_misreading.tex, 04_collectives.tex, 05_agency.tex; prop:col-whole-population added under decision 22 (H6-P1, E3 proof check of col)

The main results of \cref{sec:coh,sec:frm,sec:sca,sec:min,sec:abs} and those of the four sections that develop hypothesis H5 (\cref{sec:met,sec:hie,sec:ren,sec:env}) are listed in the same form in the supplement (\cref{tab:int-results-b,tab:int-results-c,tab:int-results-d}).

What is not established is listed in the supplement (\cref{sec:int-open}): in particular, the formation chain is a theorem only in the two-group model at infinite policy precision with exact in-group translations, at the levels of a window that is nonempty by computation, and from a declared start that is itself the split, contrast is defined only for diffusive world-conditional precisions, no population analyzed here is shown to meet every clause of \cref{def:agy-meta-agent}, and hypothesis H5a has no approximate form beyond a conditional error term (Supplement, \cref{thm:apx-posterior}). Throughout, the scope is that of the proofs: linear-Gaussian dynamics, mostly with fixed or exogenously scheduled attention; literal receivers in the network results; stationarity wherever a length or time is defined; and, for meta-agents and hierarchies, coarse states, partitions and comparison geometries declared before the analysis (\cref{asm:int-a12}). The results stated for other models are named in the same subsection of the supplement. The open problems that bear most on the five ideas are gathered in \cref{sec:dis-open}.

\Cref{sec:opn} takes up ten open problems posed in the recent literature on collective active inference. For each it states the question as posed, what the main text and the extensions of the supplement establish, and what remains open. None is described as solved: every statement holds in the model in which it is proved, no problem is closed in the generality in which it was posed, and two, energetic viability and evolutionary individuality, are addressed only by a scoped model. \Cref{tab:opn-status,tab:opn-status-b} list the status of each.

The framework is distinct from the author's gauge-theoretic manuscripts \citep{dennis2026magent,dennis2026participatory} and does not take them as a foundation: the relation is discussed in \cref{sec:dis-gauge}, and the supplement (\cref{sec:int-gauge}) lists the results imported from that program and relates the framework to the minimal model of the author's earlier report \citep{dennis2026communicating}, which \cref{sec:min} places in the framework.

\subsection{Organization}\label{sec:int-organization}

The main text proceeds as follows. \Cref{sec:fw} sets up codes, reception and the maps $T$ and $O$, interpretation levels, channels and the level-0 network, and \cref{sec:mis} develops hypothesis H2. \Cref{sec:col} treats boundaries, contrast, size and coarse media, and \cref{sec:agy} closure, agency and interaction necessity, ending with H1 in final form. \Cref{sec:coh} develops coherence, \cref{sec:frm} formation and \cref{sec:sca} the conditional theorems behind H4; \cref{sec:min} revisits the minimal model of the earlier report, and \cref{sec:abs} treats cultures, economies, nations and cities. Four sections develop hypothesis H5: meta-agents as agents (\cref{sec:met}), hierarchy (\cref{sec:hie}), renormalization (\cref{sec:ren}) and abstract environments with selective sensing (\cref{sec:env}). \Cref{sec:opn} sets the framework against ten open problems posed in the literature, \cref{sec:pre} collects predictions and what would refute them, \cref{sec:rel} relates the framework to prior work, \cref{sec:dis} discusses the status of the hypotheses and the limitations, and \cref{sec:not} lists the notation.

The supplement has four parts, and every number that carries the prefix S refers to it. The first holds, section by section, the secondary results, examples and extended discussion that the main text cites (\cref{sup:int,sup:fw,sup:mis,sup:col,sup:agy,sup:coh,sup:frm,sup:sca,sup:min,sup:abs,sup:met,sup:hie,sup:ren,sup:env,sup:pre,sup:rel,sup:dis}). The second holds eight extensions that address the open problems: approximate renormalization (\cref{sec:apx}, for \cref{sec:opn-1,sec:opn-2}), recursive formation (\cref{sec:rfm}, \cref{sec:opn-3}), conflicting objectives and scale laws (\cref{sec:obj}, \cref{sec:opn-4,sec:opn-7}), provenance (\cref{sec:prv}, \cref{sec:opn-5}), learned interfaces (\cref{sec:lrn}, \cref{sec:opn-6}), energetics (\cref{sec:nrg}, \cref{sec:opn-8}), evolutionary individuality (\cref{sec:evo}, \cref{sec:opn-9}) and identification (\cref{sec:idn}, \cref{sec:opn-10}). The third holds the proofs (\cref{app:fw,app:mis,app:col,app:agy,app:coh,app:frm,app:sca,app:min,app:abs,app:met,app:hie,app:ren,app:env,app:apx,app:rfm,app:obj,app:prv,app:lrn,app:nrg,app:evo,app:idn}), and the fourth the numerical checks (\cref{app:checks}), single-threaded double-precision confirmations on small instances.
